%% ma: L12-B08-0062
%% lop: 12
%% chuong: 2
%% bai: 8
%% dang: 12.8.35
%% loai: DS
%% muc: [NB, TH, TH, VD]
%% dap_an: "ĐĐSĐ"
%% nguon: "LaTeX giáo viên gửi 21/10/2026 (M05, Vectơ và hệ trục tọa độ trong không gian, 107 câu), Phần II Câu 19"
%% kien_thuc: "Bài 7 (tọa độ điểm), Bài 8 (biểu thức tọa độ các phép toán vectơ, trọng tâm, tích vô hướng)"
%% trang_thai: chinh-thuc
%% ghi_chu: "Giáo viên đã duyệt 21/10/2026. Nguồn ghi trong tài liệu: SGD Ninh Bình 2025-2026. Hình gốc không đúng tỉ lệ cạnh; đã vẽ lại theo AB = 5, AD = 6, AA' = 10"
\begin{ex}
Cho hình hộp chữ nhật $ABCD.A'B'C'D'$ có đáy $ABCD$ là hình chữ nhật với $AB=5$, $AD=6$, $AA'=10$; $G$ là trọng tâm của tam giác $BDA'$. Gắn hệ tọa độ $Oxyz$ có gốc $O$ trùng với điểm $A$, các tia $Ox$, $Oy$, $Oz$ lần lượt trùng với các tia $AB$, $AD$, $AA'$ (tham khảo hình bên dưới).
\begin{center}
\begin{tikzpicture}[x={(-.42cm,-.38cm)},y={(.5cm,0cm)},z={(0cm,.36cm)},scale=.85]
  \coordinate (A) at (0,0,0); \coordinate (B) at (5,0,0); \coordinate (C) at (5,6,0); \coordinate (D) at (0,6,0);
  \coordinate (A') at (0,0,10); \coordinate (B') at (5,0,10); \coordinate (C') at (5,6,10); \coordinate (D') at (0,6,10);
  \draw[khuat] (A)--(B) (A)--(D) (A)--(A');
  \draw (B)--(C)--(D)--(D')--(A')--(B')--cycle (B')--(C')--(D') (C)--(C');
  \draw[truc] (B)--(7.5,0,0) node[below]{$x$};
  \draw[truc] (D)--(0,8,0) node[right]{$y$};
  \draw[truc] (A')--(0,0,12.5) node[left]{$z$};
  \path (A) node[above right,inner sep=1pt]{$A$} (B) node[left]{$B$} (C) node[below right]{$C$} (D) node[below]{$D$}
        (A') node[left]{$A'$} (B') node[left]{$B'$} (C') node[below right]{$C'$} (D') node[above right]{$D'$};
\end{tikzpicture}
\end{center}
\choiceTF
{\True Tọa độ của điểm $C$ là $(5;6;0)$.}
{\True $\left|\overrightarrow{AC}+2\overrightarrow{A'B}\right|=\sqrt{661}$.}
{Tọa độ điểm $G$ là $\left(2;\dfrac53;\dfrac{10}3\right)$.}
{\True $\overrightarrow{AG}\cdot\overrightarrow{B'C}=-\dfrac{64}3$.}
\loigiai{$A(0;0;0)$, $B(5;0;0)$, $D(0;6;0)$, $A'(0;0;10)$, $B'(5;0;10)$.

a) $C(5;6;0)$. Đúng.

b) $\overrightarrow{AC}=(5;6;0)$, $\overrightarrow{A'B}=(5;0;-10)$, $\overrightarrow{AC}+2\overrightarrow{A'B}=(15;6;-20)$ có độ dài $\sqrt{225+36+400}=\sqrt{661}$. Đúng.

c) $G\left(\dfrac53;2;\dfrac{10}3\right)$. Sai.

d) $\overrightarrow{AG}=\left(\dfrac53;2;\dfrac{10}3\right)$, $\overrightarrow{B'C}=(0;6;-10)$, $\overrightarrow{AG}\cdot\overrightarrow{B'C}=12-\dfrac{100}3=-\dfrac{64}3$. Đúng.}
\end{ex}
